% Generated draft: requires mathematical and source review.
\documentclass[11pt]{article}
\usepackage[margin=25mm]{geometry}
\usepackage{fontspec}
\setmainfont{Latin Modern Roman}
\usepackage{amsmath,amssymb,mathtools,unicode-math}
\setmathfont{Latin Modern Math}
\usepackage{enumitem,needspace,xurl,hyperref}
\setlength{\parindent}{0pt}
\setlength{\parskip}{5pt}
\setlength{\emergencystretch}{3em}
\newcommand{\N}{\mathbb{N}}
\newcommand{\Z}{\mathbb{Z}}
\newcommand{\Q}{\mathbb{Q}}
\newcommand{\R}{\mathbb{R}}
\newcommand{\C}{\mathbb{C}}
\newcommand{\F}{\mathbb{F}}
\newcommand{\A}{\mathbb{A}}
\newcommand{\PP}{\mathbb{P}}
\newcommand{\E}{\mathbb{E}}
\newcommand{\eps}{\varepsilon}
\newcommand{\dd}{\,\mathrm{d}}
\begin{document}
\section*{UnsolvedMath ID: 1200039}
\textbf{Source problem code:} catalogue entry
\subsection{Definitions and mathematical statement}
For a finitely generated group $\Gamma$, let $d(\Gamma)$ be its minimum number of generators. Its absolute rank gradient is
\[
\operatorname{RG}(\Gamma)=
\inf_{H\leq\Gamma,\ [\Gamma:H]<\infty}
\frac{d(H)-1}{[\Gamma:H]}.
\]
For a specified nested finite-index chain $(\Gamma_i)$, a different object is the chain rank gradient
\[
\operatorname{RG}(\Gamma,(\Gamma_i))=
\lim_i\frac{d(\Gamma_i)-1}{[\Gamma:\Gamma_i]},
\]
when taken under the customary chain hypotheses; these hypotheses may require normality, trivial intersection, or a Farber condition depending on the intended question. The source does not specify a chain or which convention it uses.

A lattice $\Gamma$ in a semisimple real Lie group $S$ is a discrete subgroup such that $S/\Gamma$ has finite invariant measure. In the usual higher-real-rank setting $\operatorname{rank}_{\R}S\geq2$. The source does not spell out connectedness, finite-center, compact-factor, irreducibility, or simple-factor rank assumptions.

With the absolute convention, the question is whether
\[
\operatorname{RG}(\Gamma)=0
\quad\text{for every higher-rank semisimple real lattice }\Gamma.
\]
Equivalently, for each such $\Gamma$ and each $\varepsilon>0$, there should be a finite-index subgroup $H\leq\Gamma$ with
\[
\frac{d(H)-1}{[\Gamma:H]}<\varepsilon.
\]
If a chain convention is intended, vanishing for one chosen chain, existence of a suitable chain, and vanishing for every admissible chain are distinct quantifier choices and must not be silently identified with this formulation.

The supplied research summary reports an affirmative result for \emph{right-angled higher-rank lattices}, citing AGN (2017), and says this covers most non-uniform lattices. It leaves the assertion for every higher-rank semisimple real lattice, particularly arbitrary uniform (cocompact) lattices, as an open conjecture. ``Right-angled'' and ``most'' are not defined in the record, and the precise rank-gradient convention of the cited theorem is not supplied. No extension to all cocompact lattices is asserted here.
\subsection{Short English statement}
The problem asks whether the number of generators of finite-index subgroups grows sublinearly in their index for every higher-rank real semisimple lattice. The supplied summary reports a right-angled-lattice case, with the general case still open. The absolute-versus-chain rank-gradient convention requires clarification.
\subsection{Sources}
\begin{itemize}
\item Local UnsolvedMath registry, permanent ID 1200039, exact source code catalogue entry, ``Some Questions --- Question 39''; status partially solved, classification PARTIAL-PROGRESS.
\item Ab\'ert, \emph{Some questions} (2010), Question 39, \url{https://www.renyi.hu/~abert/questions.pdf}. The provenance records manual transcription of a public author-hosted PDF, accessed July 29, 2026; the 2010 list was inspected, with current status not independently verified.
\item ``AGN,'' \emph{Duke Mathematical Journal} 166 (2017), as given in the supplied summary for the right-angled higher-rank lattice result. The record does not expand the author initials, give a title or page range, define the right-angled condition, or state the exact gradient hypotheses. The summary's claimed scope, ``most non-uniform lattices,'' is retained as an unquantified source assessment.
\end{itemize}
\end{document}
